\documentclass[12pt]{article}

\usepackage[a4paper, top=2cm, bottom=2cm, left=2.5cm, right=2.5cm]{geometry}

\usepackage{amsmath, amssymb}
\usepackage{tikz}
\usetikzlibrary{arrows.meta,positioning,fit,calc}
\usepackage{multirow}
\usepackage{array}
\usepackage{siunitx}
\usepackage{enumitem}
\usepackage{float}
\usepackage{fontspec}
\usetikzlibrary{decorations.pathmorphing}
\usepackage{pdflscape}
\usepackage{tabularx}
\usepackage{booktabs}
\usepackage{graphicx}
\usepackage{listings}
\usepackage[table]{xcolor}
\usepackage{hyperref}
\usepackage{bookmark}
\usepackage{pdfpages}

% ---------------- FONTS ----------------
\setmainfont[
    UprightFont = Handlee-Regular.ttf,
    AutoFakeBold = 2.5
]{Handlee}

\newfontfamily\headingfont{PT.ttf}
\setmonofont{DejaVu Sans Mono}

% ---------------- HANDWRITTEN MATH ----------------
\usepackage{mathastext}
\Mathastext
\everymath{\displaystyle}

% ---------------- GENERAL ----------------
\setlength{\parindent}{0pt}
\pagenumbering{gobble}
\linespread{1.2}

\newcommand{\mychapter}[2]{%
    \phantomsection
    \hypertarget{#2}{}%
    \bookmark[level=0,dest=#2]{#1}%
    \begin{center}
        {\headingfont\Huge #1}
        \par
        \vspace{4pt}
        \hrule
    \end{center}
}

\newcommand{\mysection}[3]{%
    \phantomsection
    \hypertarget{#3}{}%
    \bookmark[level=1,dest=#3]{#1\space #2}%
    \newpage
    \noindent{\headingfont\Large #1\space #2\par}
    \vspace{4pt}
    \hrule
    \vspace{15pt}
}

\newcommand{\mysubsection}[2]{%
    \vspace{5px}
    \noindent{\headingfont\large #1\space #2\par}
    \vspace{4pt}
    \hrule
    \vspace{15pt}
}

\newcommand{\insertfigure}[2]{%
    \begin{figure}[H]
        \centering
        \fbox{\includegraphics[width=#1\textwidth]{../2. Figures/#2}}
    \end{figure}
}

\begin{document}

\mychapter{Chapter 5 : Mass and Energy Analysis of Control Volumes}{6f4y5kdh}

\begin{center}
\renewcommand{\arraystretch}{1.5}

\begin{tabularx}{0.9\textwidth}{
|>{\centering\arraybackslash}p{0.14\textwidth}
|>{\raggedright\arraybackslash}X|
}
\hline

\textbf{Section} &
\textbf{Core Revision Focus}
\\
\hline

5.1 &
Conservation of mass for control volumes; mass and volume flow rates, normal velocity, general mass balance, steady-flow mass balance, and the incompressible-flow continuity relation.
\\
\hline

5.2 &
Flow work and the energy carried by flowing fluid; \(w_{\mathrm{flow}}=Pv\), enthalpy \(h=u+Pv\), and total flowing-fluid energy \(h+\frac{V^2}{2}+gz\).
\\
\hline

5.3 &
Steady-flow energy analysis; mass and energy balances, the steady-flow energy equation, sign conventions, and identification of negligible kinetic, potential, heat-transfer, and work terms.
\\
\hline

5.4 &
Application of mass and energy balances to steady-flow engineering devices: nozzles, diffusers, turbines, compressors, pumps, fans, throttling valves, mixing chambers, heat exchangers, and pipe/duct flow.
\\
\hline

5.5 &
Unsteady-flow energy analysis; mass and energy accumulation, uniform-flow processes, distinction between flowing-fluid enthalpy and stored-fluid internal energy, and charging/discharging of tanks.
\\
\hline



\end{tabularx}
\end{center}

% ============================================================
\mysection{5.1}{Conservation of Mass}{5.1}

\mysubsection{}{Mass and Volume Flow Rates}

For flow through a control surface, only the velocity component normal to the area contributes.

\[
\delta \dot{m}=\rho V_n\,dA_c
\]

\[
\dot{m}
=
\int_{A_c}\rho V_n\,dA_c
\]

\[
V_{\mathrm{avg}}
=
\frac{1}{A_c}\int_{A_c}V_n\,dA_c
\]

For uniform/average velocity:

\[
\boxed{\dot{m}=\rho V_{\mathrm{avg}}A_c}
\]

\[
\dot{V}
=
\int_{A_c}V_n\,dA_c
=
V_{\mathrm{avg}}A_c
=
VA_c
\]

\[
\boxed{
\dot{m}
=
\rho\dot{V}
=
\frac{\dot{V}}{v}
}
\qquad
\boxed{
m=\rho V=\frac{V}{v}
}
\]

\[
\rho=\frac{1}{v}
\]

\mysubsection{}{Conservation of Mass Principle}

\[
\boxed{
\text{Mass entering}
-
\text{Mass leaving}
=
\text{Change in mass within CV}
}
\]

\[
m_{\mathrm{in}}-m_{\mathrm{out}}
=
\Delta m_{\mathrm{CV}}
\]

\[
\Delta m_{\mathrm{CV}}
=
m_{\mathrm{final}}-m_{\mathrm{initial}}
\]

Rate form:

\[
\boxed{
\dot{m}_{\mathrm{in}}
-
\dot{m}_{\mathrm{out}}
=
\frac{dm_{\mathrm{CV}}}{dt}
}
\]

Mass contained in the control volume:

\[
dm=\rho\,dV
\]

\[
m_{\mathrm{CV}}
=
\int_{\mathrm{CV}}\rho\,dV
\]

\[
\frac{dm_{\mathrm{CV}}}{dt}
=
\frac{d}{dt}
\int_{\mathrm{CV}}\rho\,dV
\]

For constant mass inside the CV:

\[
\boxed{\frac{dm_{\mathrm{CV}}}{dt}=0}
\]

\mysubsection{}{Normal Component of Velocity}

\[
V_n=V\cos\theta
\]

\[
\boxed{
V_n=\vec V\cdot\vec n
}
\]

With outward unit normal:

\[
V_n>0
\quad\text{outflow}
\qquad
V_n<0
\quad\text{inflow}
\]

\[
\delta\dot m
=
\rho V_n\,dA
=
\rho(\vec V\cdot\vec n)\,dA
\]

\[
\boxed{
\dot m_{\mathrm{net}}
=
\int_{\mathrm{CS}}
\rho(\vec V\cdot\vec n)\,dA
}
\]

\mysubsection{}{General Conservation of Mass Equation}

\[
\boxed{
\frac{dm_{\mathrm{CV}}}{dt}
+
\dot m_{\mathrm{out}}
-
\dot m_{\mathrm{in}}
=
0
}
\]

Equivalent form:

\[
\boxed{
\frac{dm_{\mathrm{CV}}}{dt}
=
\sum_{\mathrm{in}}\dot m
-
\sum_{\mathrm{out}}\dot m
}
\]

Integral form:

\[
\boxed{
\frac{d}{dt}\int_{\mathrm{CV}}\rho\,dV
=
\sum_{\mathrm{in}}\dot m
-
\sum_{\mathrm{out}}\dot m
}
\]

If velocity is normal to the surface:

\[
\vec V\cdot\vec n=V
\]

\[
\int_A\rho(\vec V\cdot\vec n)\,dA
=
\rho VA
\]

\mysubsection{}{Mass Balance for Steady-Flow Processes}

Steady flow:

\[
m_{\mathrm{CV}}=\text{constant}
\qquad\Longrightarrow\qquad
\frac{dm_{\mathrm{CV}}}{dt}=0
\]

Hence,

\[
\boxed{
\sum_{\mathrm{in}}\dot m
=
\sum_{\mathrm{out}}\dot m
}
\]

For a single inlet and outlet:

\[
\boxed{\dot m_1=\dot m_2}
\]

\[
\dot m=\rho VA
\]

\[
\boxed{
\rho_1V_1A_1
=
\rho_2V_2A_2
}
\]

\mysubsection{}{Special Case: Incompressible Flow}

For constant density:

\[
\boxed{
\sum_{\mathrm{in}}\dot V
=
\sum_{\mathrm{out}}\dot V
}
\]

For one inlet and one outlet:

\[
\boxed{
\dot V_1=\dot V_2
}
\]

\[
\boxed{
V_1A_1=V_2A_2
}
\]

For incompressible steady flow:

\[
\boxed{
\dot m_{\mathrm{in}}=\dot m_{\mathrm{out}}
}
\]

but, in general,

\[
\boxed{
\dot V_{\mathrm{in}}\neq\dot V_{\mathrm{out}}
}
\]

when density changes.

% ============================================================
\mysection{5.2}{Flow Work and the Energy of a Flowing Fluid}{5.2}

\mysubsection{}{Flow Work}

Pressure force at an inlet/exit:

\[
F=PA
\]

Flow work:

\[
W_{\mathrm{flow}}
=
FL
=
PAL
=
PV
\]

Specific flow work:

\[
\boxed{
w_{\mathrm{flow}}=Pv
}
\]

Thus the pressure/flow energy carried with each unit mass is \(Pv\).

\mysubsection{}{Total Energy of a Flowing Fluid}

Specific total energy of a nonflowing fluid:

\[
e=u+ke+pe
\]

\[
\boxed{
e=u+\frac{V^2}{2}+gz
}
\]

For flowing fluid, include flow work:

\[
\theta
=
Pv+e
\]

\[
\theta
=
Pv+u+\frac{V^2}{2}+gz
\]

Define enthalpy:

\[
\boxed{
h=u+Pv
}
\]

Therefore,

\[
\boxed{
\theta
=
h+\frac{V^2}{2}+gz
}
\]

\[
\boxed{
\text{Flowing fluid: }h
\qquad
\text{Nonflowing fluid: }u
}
\]

\mysubsection{}{Energy Transport by Mass}

Energy transported with mass:

\[
\boxed{
E_{\mathrm{mass}}=m\theta
}
\]

\[
E_{\mathrm{mass}}
=
m
\left(
h+\frac{V^2}{2}+gz
\right)
\]

Rate of energy transport:

\[
\boxed{
\dot E_{\mathrm{mass}}
=
\dot m\theta
=
\dot m
\left(
h+\frac{V^2}{2}+gz
\right)
}
\]

When kinetic and potential energies are negligible:

\[
\boxed{
E_{\mathrm{mass}}\approx mh
\qquad
\dot E_{\mathrm{mass}}\approx\dot mh
}
\]

\mysubsection{}{Energy Transport Through an Inlet or Exit}

For an infinitesimal mass:

\[
\delta E_{\mathrm{mass}}
=
\theta\,\delta m
\]

\[
E_{\mathrm{in,mass}}
=
\int_{m_i}\theta_i\,\delta m_i
\]

\[
\boxed{
E_{\mathrm{in,mass}}
=
\int_{m_i}
\left(
h_i+\frac{V_i^2}{2}+gz_i
\right)
\delta m_i
}
\]

% ============================================================
\mysection{5.3}{Energy Analysis of Steady-Flow Systems}{5.3}

\mysubsection{}{Characteristics of a Steady-Flow Process}

For a steady-flow control volume:

\[
\boxed{
V_{\mathrm{CV}}=\text{constant}
}
\]

\[
\boxed{
m_{\mathrm{CV}}=\text{constant}
}
\]

\[
\boxed{
E_{\mathrm{CV}}=\text{constant}
}
\]

Therefore, there is no accumulation of mass or energy within the CV.

\mysubsection{}{Mass Balance for Steady-Flow Systems}

\[
\boxed{
\sum_{\mathrm{in}}\dot m
=
\sum_{\mathrm{out}}\dot m
}
\]

Single stream:

\[
\dot m_1=\dot m_2=\dot m
\]

\[
\boxed{
\dot m=\rho VA
}
\]

\[
\boxed{
\rho_1V_1A_1
=
\rho_2V_2A_2
}
\]

\mysubsection{}{Energy Balance for Steady-Flow Systems}

Since \(E_{\mathrm{CV}}\) is constant,

\[
\Delta E_{\mathrm{CV}}=0
\]

Thus,

\[
\dot E_{\mathrm{in}}
=
\dot E_{\mathrm{out}}
\]

Using the sign convention of the chapter,

\[
\boxed{
\dot Q_{\mathrm{in}}
+
\dot W_{\mathrm{in}}
+
\sum_{\mathrm{in}}\dot m\theta
=
\dot Q_{\mathrm{out}}
+
\dot W_{\mathrm{out}}
+
\sum_{\mathrm{out}}\dot m\theta
}
\]

where

\[
\boxed{
\theta=h+\frac{V^2}{2}+gz
}
\]

Hence,

\[
\boxed{
\dot Q-\dot W
=
\sum_{\mathrm{out}}
\dot m
\left(
h+\frac{V^2}{2}+gz
\right)
-
\sum_{\mathrm{in}}
\dot m
\left(
h+\frac{V^2}{2}+gz
\right)
}
\]

Here,

\[
\dot Q=\dot Q_{\mathrm{net,in}}
\]

\[
\dot W=\dot W_{\mathrm{net,out}}
\]

\mysubsection{}{Single-Stream Steady-Flow Energy Equation}

For one inlet and one outlet:

\[
\dot m_1=\dot m_2=\dot m
\]

\[
\boxed{
\dot Q-\dot W
=
\dot m
\left[
h_2-h_1
+
\frac{V_2^2-V_1^2}{2}
+
g(z_2-z_1)
\right]
}
\]

Per unit mass:

\[
\boxed{
q-w
=
h_2-h_1
+
\frac{V_2^2-V_1^2}{2}
+
g(z_2-z_1)
}
\]

\[
q=\frac{\dot Q}{\dot m}
\qquad
w=\frac{\dot W}{\dot m}
\]

\mysubsection{}{Negligible Kinetic and Potential Energy Changes}

If

\[
\Delta ke\approx0
\qquad
\Delta pe\approx0
\]

then

\[
\boxed{
q-w=h_2-h_1
}
\]

\mysubsection{}{Heat Transfer}

\[
\dot Q>0
\quad\text{heat enters CV}
\]

\[
\dot Q<0
\quad\text{heat leaves CV}
\]

\[
\dot Q=0
\quad\text{adiabatic}
\]

\mysubsection{}{Work Transfer}

For no work interaction:

\[
\boxed{\dot W=0}
\]

\mysubsection{}{Enthalpy Change}

\[
\boxed{
\Delta h=h_2-h_1
}
\]

For an ideal-gas-like constant-\(c_p\) approximation:

\[
\boxed{
\Delta h
=
c_{p,\mathrm{avg}}(T_2-T_1)
}
\]

Unit relation:

\[
(\mathrm{kg/s})(\mathrm{kJ/kg})
=
\mathrm{kW}
\]

\mysubsection{}{Kinetic Energy Change}

\[
\boxed{
\Delta ke
=
\frac{V_2^2-V_1^2}{2}
}
\]

\[
\mathrm{m^2/s^2}
=
\mathrm{J/kg}
\]

If \(V_1\approx V_2\):

\[
\boxed{
\Delta ke\approx0
}
\]

\mysubsection{}{Potential Energy Change}

\[
\boxed{
\Delta pe=g(z_2-z_1)
}
\]

\mysubsection{}{Steady-Flow Energy Equation: Summary}

\[
\boxed{
\dot Q-\dot W
=
\sum_{\mathrm{out}}
\dot m
\left(
h+\frac{V^2}{2}+gz
\right)
-
\sum_{\mathrm{in}}
\dot m
\left(
h+\frac{V^2}{2}+gz
\right)
}
\]

Single stream:

\[
\boxed{
\dot Q-\dot W
=
\dot m
\left[
h_2-h_1
+
\frac{V_2^2-V_1^2}{2}
+
g(z_2-z_1)
\right]
}
\]

Per unit mass:

\[
\boxed{
q-w
=
h_2-h_1
+
\frac{V_2^2-V_1^2}{2}
+
g(z_2-z_1)
}
\]

If \(\Delta ke,\Delta pe\approx0\):

\[
\boxed{
q-w=h_2-h_1
}
\]

\mysection{5.4}{Some Steady-Flow Engineering Devices}{5.4}
% ============================================================
% Table 1: Nozzles, Diffusers, Turbines, Compressors, Pumps, Fans
% ============================================================

\renewcommand{\arraystretch}{1.5}

\begin{center}
\begin{tabularx}{\textwidth}{
|>{\bfseries\raggedright\arraybackslash}p{0.15\textwidth}
|>{\raggedright\arraybackslash}p{0.22\textwidth}
|>{\raggedright\arraybackslash}p{0.35\textwidth}
|>{\raggedright\arraybackslash}X|
}
\hline
\textbf{Device Name} &
\textbf{Function} &
\textbf{Governing Equation} &
\textbf{Assumption} \\
\hline

Nozzle &
Converts enthalpy into kinetic energy; velocity increases. &
$\displaystyle
h_1+\frac{V_1^2}{2}
=
h_2+\frac{V_2^2}{2}$ &
$\displaystyle
\dot Q\approx0,\quad
\dot W=0,\quad
\Delta pe\approx0$ \\
\hline

Diffuser &
Converts kinetic energy into enthalpy; velocity decreases. &
$\displaystyle
h_1+\frac{V_1^2}{2}
=
h_2+\frac{V_2^2}{2}$ &
$\displaystyle
\dot Q\approx0,\quad
\dot W=0,\quad
\Delta pe\approx0$ \\
\hline

Turbine &
Produces shaft work from the flowing fluid. &
$\displaystyle
w_{\mathrm{out}}=h_1-h_2$ &
$\displaystyle
\dot Q\approx0,\quad
\Delta ke\approx0,\quad
\Delta pe\approx0$ \\
\hline

Compressor &
Raises the pressure of a gas by consuming shaft work. &
$\displaystyle
w_{\mathrm{in}}=h_2-h_1-q
$\newline
Adiabatic:
$\displaystyle
w_{\mathrm{in}}=h_2-h_1$ &
$\displaystyle
\dot Q\approx0,\quad
\Delta ke\approx0,\quad
\Delta pe\approx0$ \\
\hline

Pump &
Raises the pressure of a liquid by consuming shaft work. &
$\displaystyle
w_{\mathrm{in}}=h_2-h_1-q
$\newline
Adiabatic:
$\displaystyle
w_{\mathrm{in}}=h_2-h_1$ &
$\displaystyle
\dot Q\approx0,\quad
\Delta ke\approx0,\quad
\Delta pe\approx0$ \\
\hline

Fan &
Produces a relatively small pressure increase in a gas. &
$\displaystyle
w_{\mathrm{in}}=h_2-h_1-q$ &
$\displaystyle
\dot Q\approx0,\quad
\Delta ke\approx0,\quad
\Delta pe\approx0$ \\
\hline

\end{tabularx}
\end{center}


% ============================================================
% Table 2: Throttling Valve, Mixing Chamber, Heat Exchanger, Pipe/Duct
% ============================================================

\renewcommand{\arraystretch}{1.5}

\begin{center}
\begin{tabularx}{\textwidth}{
|>{\bfseries\raggedright\arraybackslash}p{0.15\textwidth}
|>{\raggedright\arraybackslash}p{0.22\textwidth}
|>{\raggedright\arraybackslash}p{0.35\textwidth}
|>{\raggedright\arraybackslash}X|
}
\hline
\textbf{Device Name} &
\textbf{Function} &
\textbf{Governing Equation} &
\textbf{Assumption} \\
\hline

Throttling Valve &
Produces a large pressure drop with no useful work output; isenthalpic device. &
$\displaystyle
\boxed{h_2\approx h_1}
$
\newline
$\displaystyle
u_1+P_1v_1=u_2+P_2v_2
$ &
$\displaystyle
q\approx0,\quad
w=0,\quad
\Delta ke\approx0,\quad
\Delta pe\approx0$ \\
\hline

Mixing Chamber &
Mixes two or more streams; mass and energy are conserved. &
$\displaystyle
\sum_{\mathrm{in}}\dot m
=
\sum_{\mathrm{out}}\dot m
$
\newline
$\displaystyle
\sum_{\mathrm{in}}\dot m h_{\mathrm{in}}
=
\sum_{\mathrm{out}}\dot m h_{\mathrm{out}}
$ &
$\displaystyle
\dot Q\approx0,\quad
\dot W=0,\quad
\Delta ke\approx0,\quad
\Delta pe\approx0$ \\
\hline

Heat Exchanger &
Transfers heat between streams without mixing them. &
$\displaystyle
\sum_{\mathrm{in}}\dot m h
=
\sum_{\mathrm{out}}\dot m h
$
\newline
$\displaystyle
\dot Q_{\mathrm{gained}}
=
\dot Q_{\mathrm{lost}}
$ &
$\displaystyle
\dot W=0,\quad
\Delta ke\approx0,\quad
\Delta pe\approx0,\quad
\dot Q_{\mathrm{CV}}\approx0$ \\
\hline

Pipe / Duct &
Transports fluid; heat transfer and work interactions may occur. &
$\displaystyle
\dot Q-\dot W
=
\dot m(h_2-h_1)
$
\newline
$\displaystyle
\dot Q-\dot W
=
\dot m c_p(T_2-T_1)
$ &
Steady flow,
$\displaystyle
\Delta ke\approx0,\quad
\Delta pe\approx0$ \\
\hline

\end{tabularx}
\end{center}

% ============================================================
\mysection{5.5}{Energy Analysis of Unsteady-Flow Processes}{5.5}

\mysubsection{}{Characteristics of Unsteady-Flow Systems}

\textbf{\ Steady flow:}

\[
\boxed{
\frac{dm_{\mathrm{CV}}}{dt}=0,
\qquad
\frac{dE_{\mathrm{CV}}}{dt}=0
}
\]

\textbf{\ Unsteady flow:}

Mass and/or energy within the control volume may change with time.

The CV is usually stationary, but its boundary may move; therefore boundary work can occur.

\[
\boxed{
\text{Unsteady flow}\neq\text{steady flow}
}
\]

\mysubsection{}{Mass Balance for an Unsteady-Flow Process}

\[
\boxed{
m_{\mathrm{in}}-m_{\mathrm{out}}
=
\Delta m_{\mathrm{system}}
}
\]

\[
\Delta m_{\mathrm{system}}
=
m_{\mathrm{final}}-m_{\mathrm{initial}}
\]

Thus,

\[
\boxed{
m_{\mathrm{in}}-m_{\mathrm{out}}
=
m_{\mathrm{final}}-m_{\mathrm{initial}}
}
\]

For a control volume:

\[
\boxed{
m_i-m_e
=
(m_2-m_1)_{\mathrm{CV}}
}
\]

Useful special cases:

\[
m_i=0
\quad\text{no inlet}
\]

\[
m_e=0
\quad\text{no exit}
\]

\[
m_1=0
\quad\text{initially evacuated CV}
\]

\mysubsection{}{Energy Changes in an Unsteady-Flow Process}

\[
\boxed{
E_{\mathrm{in}}-E_{\mathrm{out}}
=
\Delta E_{\mathrm{system}}
}
\]

\[
\boxed{
\text{Net energy entering}
=
\text{Change in energy stored}
}
\]

Energy may cross the boundary through:

\[
\boxed{
\text{heat}+\text{work}+\text{mass flow}
}
\]

\mysubsection{}{Uniform-Flow Process}

A uniform-flow process is an idealization used for practical unsteady-flow processes.

The inlet/exit properties are treated as uniform and steady during the process, although the control volume itself may change state with time.

\[
\boxed{
\text{Uniform flow}\neq\text{steady flow}
}
\]

The state of the exiting mass at an instant is taken to correspond to the state of fluid within the CV at that instant.

\mysubsection{}{Energy of Flowing and Nonflowing Fluid}

For flowing fluid:

\[
\boxed{
\theta
=
h+ke+pe
}
\]

\[
ke=\frac{V^2}{2}
\qquad
pe=gz
\]

Therefore,

\[
\boxed{
\theta
=
h+\frac{V^2}{2}+gz
}
\]

For nonflowing fluid within the CV:

\[
\boxed{
e
=
u+ke+pe
}
\]

\[
\boxed{
e
=
u+\frac{V^2}{2}+gz
}
\]

Key distinction:

\[
\boxed{
\text{Flowing fluid}\rightarrow h
\qquad
\text{Nonflowing fluid}\rightarrow u
}
\]

\mysubsection{}{General Energy Equation for a Uniform-Flow Process}

\[
\boxed{
\left(
Q_{\mathrm{in}}
+
W_{\mathrm{in}}
+
\sum_{\mathrm{in}}m\theta
\right)
-
\left(
Q_{\mathrm{out}}
+
W_{\mathrm{out}}
+
\sum_{\mathrm{out}}m\theta
\right)
=
(m_2e_2-m_1e_1)_{\mathrm{system}}
}
\]

with

\[
\theta=h+\frac{V^2}{2}+gz
\]

\[
e=u+\frac{V^2}{2}+gz
\]

\mysubsection{}{Negligible Kinetic and Potential Energy Changes}

If

\[
\Delta ke\approx0,
\qquad
\Delta pe\approx0
\]

then

\[
\boxed{
Q-W
=
\sum_{\mathrm{out}}mh
-
\sum_{\mathrm{in}}mh
+
(m_2u_2-m_1u_1)_{\mathrm{system}}
}
\]

where

\[
Q=Q_{\mathrm{net,in}}
=
Q_{\mathrm{in}}-Q_{\mathrm{out}}
\]

\[
W=W_{\mathrm{net,out}}
=
W_{\mathrm{out}}-W_{\mathrm{in}}
\]

\mysubsection{}{Physical Meaning of the Simplified Energy Equation}

\[
\boxed{
Q-W
=
\sum_{\mathrm{out}}mh
-
\sum_{\mathrm{in}}mh
+
(m_2u_2-m_1u_1)_{\mathrm{system}}
}
\]

Unlike a steady-flow device, an unsteady-flow CV can store or release energy as its internal state changes.

\[
\boxed{
\text{Energy in}
=
\text{energy out}
+
\text{energy accumulation}
}
\]


\mysubsection{}{Charging of a Rigid Tank}

For one inlet, no outlet:

\[
m_i=m_2-m_1
\]

Initially evacuated:

\[
m_1=0
\qquad\Longrightarrow\qquad
m_i=m_2
\]

For a rigid, insulated tank with no work and negligible kinetic/potential-energy changes:

\[
Q=0,
\qquad
W=0
\]

Energy balance:

\[
m_i h_i
=
m_2u_2-m_1u_1
\]

For an initially evacuated tank:

\[
m_i h_i=m_2u_2
\]

Using

\[
m_i=m_2
\]

gives

\[
\boxed{
u_2=h_i
}
\]

Since

\[
h=u+Pv
\]

the flow energy \(Pv\) entering with the fluid can become internal energy after the fluid comes to rest inside the tank.

\mysubsection{}{Alternative Closed-System Interpretation of Tank Charging}

Boundary work for the imaginary closed-system representation:

\[
\boxed{
W_{b,\mathrm{in}}
=
-\int_1^2P_i\,dV
}
\]

At constant pressure:

\[
\boxed{
W_{b,\mathrm{in}}
=
-P_i(V_2-V_1)
}
\]

The resulting energy relation is

\[
u_2=u_i+P_iv_i
\]

Since

\[
h_i=u_i+P_iv_i
\]

\[
\boxed{
u_2=h_i
}
\]

Thus the closed-system and uniform-flow interpretations give the same result.

\mysubsection{}{Discharging of a Tank}

For one exit, no inlet:

\[
\boxed{
m_e=m_1-m_2
}
\]

For a rigid tank:

\[
W_b=0
\]

For an insulated tank:

\[
Q\approx0
\]

With negligible kinetic and potential-energy changes:

\[
\boxed{
W_{\mathrm{in}}
-
m_eh_e
=
m_2u_2-m_1u_1
}
\]

Electrical work supplied to the tank can compensate for the energy carried away by the discharged fluid.

\mysubsection{}{Constant-Temperature Discharge}

For constant temperature:

\[
T_2=T_1
\]

For a fluid for which \(u=u(T)\):

\[
\boxed{
u_2=u_1
}
\]

The process may have no accumulation at a particular state while mass and energy inside the CV still vary with time:

\[
\boxed{
\text{No accumulation within CV}
}
\]

\[
\boxed{
\text{Mass and energy within CV may change with time}
}
\]

The essential unsteady-flow interactions are:

\[
\boxed{
\text{Mass transfer}
+
\text{Energy transfer}
+
\text{Energy storage}
}
\]


\end{document}